\section{Introduction}
\label{sec:intro}

\noindent Consider a consortium of banks that together want to catch money
laundering. Fraud is fundamentally \emph{relational}: illicit funds move through
chains of accounts, shell companies, and intermediaries, and it is precisely the
\emph{structure} of these money flows---not any single transaction in
isolation---that betrays criminal intent. Graph Neural Networks (GNNs) have
therefore become the tool of choice for financial crime detection, credit
scoring, and other high-stakes risk assessments, because their message-passing
mechanism learns directly from the web of relationships between users, accounts,
and transactions \cite{kipf2016semi,velickovic2017graph,dou2020enhancing}. Yet
the very banks that would benefit most from a shared model are legally forbidden
from pooling their customers' data: privacy regulations such as the GDPR, and
plain commercial confidentiality, keep each institution's transaction graph
locked inside its own walls. Federated Learning (FL) offers a way out. By
exchanging model updates instead of raw data, a group of institutions can train
a single, more powerful detector without any of them surrendering its records
\cite{mcmahan2017communication}.

This picture, however, is dangerously incomplete. A model that is merely
\emph{accurate} is no longer fit to be deployed in decisions that shape people's
access to credit, liberty, or financial services. Three converging forces---the
European Commission's Ethics Guidelines for Trustworthy AI, the EU AI Act, and
the U.S.\ NIST AI Risk Management Framework---now demand that automated
decision systems be simultaneously \emph{fair}, \emph{privacy-preserving},
\emph{robust}, \emph{transparent}, and \emph{accountable}. These are not
optional extras: for a high-risk system under the EU AI Act they are legal
prerequisites for going to market. Meeting them is hard even for a centralised
model. In the federated, graph-structured setting they collide, and each one
becomes markedly harder than in isolation.

\paragraph{Why the federated graph setting is uniquely hard.}
The difficulty is not that we must bolt four independent features onto a
classifier; it is that the natural solution to each requirement actively
sabotages the others. \emph{Fairness} is confounded by the graph itself:
homophilous message passing means that a member of a demographic group who
happens to be connected to risky neighbours will be flagged as risky too, so the
model learns to launder a protected attribute through graph structure---and the
central server, which never sees any data, cannot even measure this bias, let
alone correct it. \emph{Privacy} appears to have a textbook answer,
differentially private SGD, but that answer fails badly here: adding calibrated
noise to every coordinate of the gradient injects perturbation whose magnitude
grows with the square root of the model dimension, and on the compact GNNs used
for tabular financial graphs this noise simply drowns the signal, collapsing
accuracy to little better than chance (we quantify this in
Section~\ref{sec:exp-privacy}). \emph{Robustness} is threatened not only by the
classical malfunctioning or malicious client, but---as we show---by a subtler
adversary that the fairness machinery itself invites: a client can transmit a
deliberately biased model while \emph{reporting} exemplary fairness statistics,
thereby hijacking any aggregator that rewards fairness. And \emph{transparency}
is undermined by the opacity of relational models, exactly when regulation
requires that each adverse decision be explainable and auditable.

\paragraph{The gap.}
The research literature has, understandably, attacked these problems one at a
time. Fair-GNN methods deliver impressive demographic parity but assume all data
sits in one place \cite{dai2021say,yang2024fairsin,lee2025dabgnn}. Fair federated
learning has matured for tabular and image data but rarely touches graphs
\cite{ezzeldin2023fairfed,li2020fair}. Byzantine-robust aggregation hardens FL
against poisoning but is oblivious to fairness \cite{blanchard2017krum,
yin2018byzantine}. Differentially private FL protects data but, as noted, at a
utility cost that is prohibitive for graph models. To the best of our knowledge,
\emph{no} existing method delivers fairness, differential privacy, and Byzantine
robustness together for federated GNNs, and none has been validated on a
large-scale cryptocurrency graph---arguably the most consequential modern fraud
domain. Closing this gap is the goal of this paper.

\paragraph{Our approach.}
We present \textbf{FedFairGNN}, a single framework in which fairness, privacy,
robustness, and transparency are designed to reinforce rather than undermine one
another. Its architecture (Fig.~\ref{fig:architecture}) intervenes at the three
natural control points of federated graph training---inside message passing, at
the local gradient, and at server aggregation---and wraps the result in an
instrumentation layer that measures, explains, and audits trustworthiness. The
key technical insight, and the reason privacy becomes almost free, is a change
of perspective on \emph{where} the sensitive attribute actually influences
learning: it does so only through a two-dimensional summary statistic, so we can
protect that statistic directly instead of blindly noising the entire gradient.

\medskip
\noindent Concretely, our contributions are:
\begin{enumerate}
\item \textbf{Fairness-Sensitive Edge Reweighting (FSER)}, a learnable attention
correction that detects and suppresses the biased cross-group edges responsible
for structural discrimination, mitigating unfairness at its source inside the
GNN rather than patching it after the fact (Section~\ref{sec:fser}).

\item \textbf{Fairness-Targeted Gradient Decomposition (FTGD)}. We observe that
the sensitive attribute enters the model \emph{only} through the
demographic-parity statistics, and exploit this to obtain
$(\epsilon,\delta)$-differential privacy by privatising two bounded-sensitivity
scalars instead of the full gradient. This makes strong privacy essentially
free---precisely where standard DP-FedAvg collapses utility---and we prove the
guarantee with a self-contained R\'enyi-DP accountant (Section~\ref{sec:ftgd}).

\item \textbf{Bi-objective Frank--Wolfe Aggregation (BFWA)} and a
\textbf{Byzantine-robust variant}. The server solves a constrained optimisation
that enforces a hard, operator-chosen fairness budget during aggregation, while
distance-based screening neutralises poisoning---including the
fairness-poisoning attack of \cite{kasyap2025attack}, which we reproduce and
defend against (Sections~\ref{sec:bfwa},~\ref{sec:exp-robust}).

\item \textbf{A trustworthiness instrumentation layer} that turns abstract
principles into measured, auditable artifacts: predictive uncertainty and
calibration, GNN attention explanations with fairness attribution, a transparent
composite \emph{trust score}, sustainability accounting, and an explicit
EU~AI~Act / NIST~RMF compliance mapping with an auto-generated model card
(Sections~\ref{sec:method-trust},~\ref{sec:trust}).

\item \textbf{A comprehensive, fully reproducible evaluation} on five real
datasets spanning credit, recidivism, social, and---uniquely---large-scale
crypto (Elliptic Bitcoin, $203{,}769$ transactions) domains, against twelve
baselines that include the most recent 2025--2026 fairness-aware federated and
graph methods. FedFairGNN attains the best overall trustworthiness trade-off,
and every number and figure in this paper is regenerated from logged runs by a
single script.
\end{enumerate}

By lowering the barrier to trustworthy, infrastructure-grade collaborative AI,
this work contributes to UN Sustainable Development Goal~9 (industry,
innovation, and resilient infrastructure). The remainder of the paper is
organised as follows. Section~\ref{sec:related} situates our work within the
fairness, privacy, and robustness literatures; Section~\ref{sec:prelim}
formalises the federated graph setting and the trustworthiness objectives;
Section~\ref{sec:method} develops the FedFairGNN framework component by
component; Section~\ref{sec:experiments} reports the empirical study;
Section~\ref{sec:trust} analyses trustworthiness holistically and maps the
system to regulation; and Section~\ref{sec:discussion} discusses implications,
limitations, and broader impact before we conclude.
